\section{Mathematical formulation}
\label{sec:formulation}

\subsection{Sets, decisions and parameters}
Let $\mathcal{T}=\{0,\dots,T-1\}$ be time steps of length $\Delta t$ (h) and $\mathcal{S}$ the scenarios, each with
probability $\pi_s$. First-stage decisions $x=(x_{pv},x_{w},x_{bp},x_{be})$ are the installed PV and wind
capacity (MW), battery power (MW) and battery energy (MWh). Each satisfies
$x_k^0\le x_k\le x_k^0+\bar x_k$, where $x_k^0$ is the existing fleet and $\bar x_k$ the maximum addition. Import
capacity $x_{imp}$ (MW) and maximum reservoir volume $x_{vmax}$ (hm$^3$) are modelled as fixed first-stage
variables ($x_k=x_k^0$). This lets their reduced costs be read as marginal values (Section~\ref{sec:marginal}).
Second-stage decisions for each $(s,t)$ are the PV, wind and run-of-river generation $g^{pv},g^{w},g^{ror}$, the
reservoir hydro output $h$, spill $\sigma$ (hm$^3$/h), reservoir volume $V$ (hm$^3$), battery charge $c$,
discharge $d$ and state of charge $e$, imports $m$, exports $y$ and unserved energy $u$. All are non-negative.

\subsection{Objective}
\begin{equation}
\min\; w\sum_{k\in\{pv,w,bp,be\}} A_k\,(x_k-x_k^0)\;+\;\sum_{s}\pi_s L_s\;+\;\lambda\Big[\zeta+\frac{1}{1-\alpha}\sum_s\pi_s\phi_s\Big],
\label{eq:obj}
\end{equation}
\begin{equation}
L_s=\frac{1}{1000}\sum_{t}\Delta t\,\Big(p^m_{s,t}m_{s,t}-\rho\,p^m_{s,t}y_{s,t}+\mathrm{VOLL}\,u_{s,t}+c^h h_{s,t}+c^{\deg}d_{s,t}\Big)\quad\text{(kEUR)},
\label{eq:opex}
\end{equation}
\begin{equation}
\phi_s\ge L_s-\zeta,\qquad \phi_s\ge0,\qquad \zeta\in\R .
\label{eq:cvar}
\end{equation}
Here $A_k=C_k[\,r(1+r)^{n_k}/((1+r)^{n_k}-1)+f_k\,]$ is the annualised cost per unit, combining the capital
recovery factor with fixed O\&M as a share $f_k$ of overnight cost $C_k$. The factor $w=T\Delta t/8760$ scales
annual costs to the modelled horizon. Existing capacity is sunk and therefore not charged. The remaining
parameters are the import price $p^m$, the export-to-import price ratio $\rho$, hydro variable O\&M $c^h$ and a
battery throughput degradation proxy $c^{\deg}$. At the optimum, $\zeta$ equals the value at risk $\mathrm{VaR}_\alpha$
of the operating-cost distribution, and the bracketed term equals $\mathrm{CVaR}_\alpha$
\cite{rockafellar2000cvar,rockafellar2002cvar}. With $\lambda=0$ the model is risk neutral; with $\lambda>0$ it
penalises the expected cost of the worst $(1-\alpha)$ tail. Investment cost is identical across scenarios, so
the risk term acts on operating cost.

\subsection{Constraints}
\paragraph{Energy balance.} For every $(s,t)$, whose dual is the system marginal cost:
\begin{equation}
g^{pv}_{s,t}+g^{w}_{s,t}+g^{ror}_{s,t}+h_{s,t}+d_{s,t}+m_{s,t}+u_{s,t}=D_{s,t}+c_{s,t}+y_{s,t}.
\label{eq:balance}
\end{equation}
Curtailment is implicit: $g^{pv}_{s,t}\le f^{pv}_t x_{pv}$ and $g^{w}_{s,t}\le f^{w}_t x_{w}$, so curtailed energy is
$f^{pv}_t x_{pv}-g^{pv}_{s,t}$ and does not appear in~\eqref{eq:balance}. This avoids double counting.
Run-of-river generation satisfies $g^{ror}_{s,t}\le R_{s,t}$, where $R_{s,t}$ is its availability.

\paragraph{Reservoir.} With the energy equivalent $\kappa=\rho_w g H\eta\cdot10^6/3.6\times10^9$ (MWh/hm$^3$),
\begin{align}
V_{s,t+1}&=V_{s,t}+\Delta t\,\big(I_{s,t}-h_{s,t}/\kappa-\sigma_{s,t}\big), &
V^{\min}&\le V_{s,t}\le x_{vmax}, \label{eq:water}\\
0&\le h_{s,t}\le P^h, & h_{s,t}/\kappa+\sigma_{s,t}&\ge r^{\min}, \label{eq:hydro}\\
V_{s,0}&=V^0_s, & V_{s,T}&\ge V^T_s. \label{eq:terminal}
\end{align}
With $H=230$~m and $\eta=0.88$, $\kappa=551.6$~MWh/hm$^3$. The fixed-head assumption makes the model linear; a
variable-head model would be bilinear in $V$ and $q$ and is left for future work. In annual runs
$V^T_s=V^0_s$, which closes the seasonal cycle. For shorter windows, $V^0_s$ and $V^T_s$ come from a coarse
(6-hourly) full-year run of the same scenario. We call this the \emph{nested boundary}. It prevents an isolated
week from either exhausting or hoarding seasonal storage.

\paragraph{Battery.}
\begin{align}
e_{s,t+1}&=e_{s,t}+\Delta t\big(\eta^c c_{s,t}-d_{s,t}/\eta^d\big), &
\varepsilon\,x_{be}&\le e_{s,t}\le x_{be}, \label{eq:soc}\\
0&\le c_{s,t},\,d_{s,t}\le x_{bp}, & e_{s,0}&=s_0x_{be},\;\; e_{s,T}\ge e_{s,0},\qquad x_{be}\le\delta^{\max}x_{bp}. \label{eq:batt}
\end{align}
Simultaneous charging and discharging is not excluded by binary variables. Because $\eta^c\eta^d<1$ and
$c^{\deg}>0$, the LP has no incentive to do it unless it must dissipate energy, and the quantity is reported
ex post for every run.

\paragraph{Interconnection, shortage and budget.}
\begin{equation}
0\le m_{s,t}\le a_s\,x_{imp},\qquad 0\le y_{s,t}\le Y,\qquad 0\le u_{s,t}\le D_{s,t},\qquad
\sum_k C_k(x_k-x_k^0)\le B .
\label{eq:imports}
\end{equation}
The import limit is an aggregate net transfer capacity scaled by scenario availability $a_s$. It is not a
power-flow model.

\subsection{Representative weeks for two-stage planning}
Hourly resolution matters for batteries, and seasonal coupling matters for the reservoir
\cite{poncelet2016temporal,kotzur2018seasonal}. To retain both, investment and stochastic problems use eight
hourly weeks per scenario, starting on days 7, 52, \dots, 325 (a systematic sample about every 45.6~days). Each
(scenario, week) pair is a sub-problem with probability $\pi_s/8$. It starts at the volume $V^0_{s,w}$ of a coarse
(6-hourly) full-year run of scenario $s$ with the existing fleet. Its end volume is not fixed. Instead it is
valued at the marginal water value $\omega_{s,w}$ of that run, the negative dual of the water-balance
constraint~\eqref{eq:water} at the end of the week (kEUR/hm$^3$):
\begin{equation}
L_{s,w} \leftarrow L_{s,w} + \omega_{s,w}\,\big(\bar V^T_{s,w} - V_{s,w,T}\big),
\label{eq:watervalue}
\end{equation}
where $\bar V^T_{s,w}$ is the coarse run's end-of-week volume. This is a first-order (one-cut) future-cost
function, in the spirit of the water values used in stochastic dual dynamic programming \cite{pereira1991sddp}.
Water used beyond the coarse plan is charged at its opportunity cost, and water saved is credited. The battery is
cyclic within each week. CVaR excess variables are defined per scenario, not per week:
$\phi_s\ge\sum_{w}\tfrac{1}{8}L_{s,w}-\zeta$. The boundaries do not depend on the plan, so the stochastic
programme, the deterministic plan's evaluation, the wait-and-see problems and all frontier evaluations share
identical seasonal information. This guarantees a consistent comparison. Section~\ref{sec:chrono} compares
this approximation with full-year chronological evaluation. It also compares it with an earlier hard-boundary
variant ($V_{s,w,T}\ge\bar V^T_{s,w}$), which we found created spurious scarcity.

\subsection{Marginal value of flexibility}
\label{sec:marginal}
For a capacity $x_k$ fixed by bounds, HiGHS returns the reduced cost $\bar d_k=c_k-(A^\top y)_k$. The marginal
operating value is $\mathrm{MV}_k=-\partial\mathbb{E}[L]/\partial x_k=c_k-\bar d_k-(A_{\text{1st}}^\top y)_k$.
The last term removes contributions from first-stage rows such as the budget and duration constraints. The
optimal operating cost is a convex, piecewise-linear function of $x_k$ \cite{bertsimas1997lp}. Hence
$\mathrm{MV}_k$ is a subgradient: it lies between the finite-difference value of adding the next increment,
$\Delta^+_k$, and that of removing the last one, $\Delta^-_k$. Duals are therefore local and may be non-unique
under degeneracy. At zero capacity (for example $x_{bp}=0$, or $x_{be}>0$ with $x_{bp}=0$) the corner is
degenerate and the dual is uninformative. We therefore report finite differences alongside duals and flag
corner cases.
